L'acide benzoïque est utilisé dans la préparation des esters odorants comme le
benzoate de méthyle $\ce{C6H5-COO-CH3}$, qui est préparé a partir de la
réaction d'estérification entre l'acide benzoïque et le méthanol en présence
d'acide sulfurique selon l'équation~:
\[
    \ce{C_6H_5-COOH + CH_3-OH <=> C_6H_5-COO-CH_3 + H_2O}
\]

On réalise l'estérification à partir d'un mélange équimolaire contenant
$n=\qty{0,3}{\mole}$ d'acide benzoïque et $n=\qty{0,3}{\mole}$ de méthanol. La
constante d'équilibre $\mathrm{K}$ associée à l'équation de la réaction
d'estérification est $\mathrm{K}=4$. 
\begin{enumerate}
    \item Citer le rôle joué par l'acide sulfurique au cours de cette réaction.
    \item Dresser le tableau d'avancement correspondant à cette réaction
          d'estérification.
    \item Montrer que l'expression de $x_{\eq}$ l'avancement de la réaction à
          l'équilibre s'écrit~:
          $x_{\eq}=\frac{n\cdot\sqrt{\mathrm{K}}}{(1+\sqrt{\mathrm{K}})}$.
    \item Déterminer la composition du mélange à l'état d'équilibre du système
          chimique.
    \item Calculer la valeur du rendement $r$ de la réaction.
    \item On ajoute une quantité d'acide benzoïque au système chimique en état
          d'équilibre.

          Répondre par Vrai ou Faux aux propositions a, b et c suivantes~:
          \begin{table}[H]
              \centering
              \renewcommand{\arraystretch}{1.5}
              \begin{tabularx}{0.8\textwidth}{|
                    >{\arraybackslash\centering}p{6mm}|L|}
                  \hline a & L'équilibre du système chimique se déplace dans le
                             sens direct \\
                  \hline b & Le rendement de cette réaction augmente \\
                  \hline c & La valeur de la constante d'équilibre $\mathrm{K}$
                             augmente \\
                  \hline
              \end{tabularx}
          \end{table}
\end{enumerate}

